\documentclass{article}
\usepackage{xcolor}
\usepackage{mdframed}
\usepackage{enumitem}
\usepackage{hyperref} %Link%
\usepackage{graphicx}
\title{Panduan Pemrograman Menggunakan LibGdx}
\author{ZVKT FaSHioN}
\date{October 2026}


\definecolor{info_color}{RGB}{23,5,120}
\definecolor{green_color}{RGB}{89,120,20}
\definecolor{linkblue}{RGB}{0,100,180}

\hypersetup{
    colorlinks=true,
    urlcolor=linkblue,
    linkcolor=black,
    citecolor=black
}

\newenvironment{info}
{\begin{mdframed}[
    linewidth=2pt,
    linecolor=info_color,
    backgroundcolor=white
]
}
{\end{mdframed}}

\newenvironment{intermeso}[1]
{\begin{mdframed}[
    linewidth=2pt,
    linecolor=green_color,
    backgroundcolor=green_color!10
]
\textbf{\large Intermeso #1}\par\vspace{0pt}
}
{\end{mdframed}}

\begin{document}
\begin{titlepage}
\maketitle
\centering
\includegraphics[width=1\textwidth]{out.png}
Proses belajar adalah hal yang paling menyenangkan saat menciptakan karya.
\end{titlepage}

\section{Abstract}
\subsection{Pendahuluan}
Panduan ini adalah dokumen yang dibuat dalam menyiapkan dan pengorganisasian kode ysng diharapkan akan mempermudah dalam kerja sama tim pada pembuatan gim menggunakan LibGDX, guna mengikuti kompetisi Arkavidia 11.0 pada bidang Gamevidia, yakni sebuah kompetisi yang diselenggarakan oleh Institut Teknologi Bandung (ITB). Dengan membawa nama tim ZVKT FaSHioN.
\setlength{\parskip}{8pt}

\subsection{Tim Developer}
\begin{enumerate}[itemsep=0pt]
    \item Rio Sakti \href{https://xn--om2b.fyi}{"refmaven"} Mahendra  : Lead Developer, Game Design, Programmers, Music and SFX Producer, Story Writer, Tester. Artist. 
    \item Sirulla Humaidi : Game Design, World Design, UI/UX, Programmers, Music Composer, Story Writer, Post Production.
    \item Muhammad Pasha Ali Rasyid : World Design, Programmers, Artist, Tester.
    \item Raffaele Hanry Aditiya : Programmer, Music COmposer, Artist, Tester.
    \item M. Raihan Pratama : And so on.
\end{enumerate}

\subsection{Garis Besar Dokumen}
Dokumen ini berisikan hal berikut
\begin{enumerate}[itemsep=0pt]
    \item Setup dan instalasi LibGDX
    \item Struktur proyek dan organisasi kode
    \item Pembagian tugas antaranggota tim
\end{enumerate}

\newpage
\tableofcontents
\newpage

\section{Java}
\subsection{Apa itu Java}

\begin{intermeso}{Pernah Denger tentang Java??}
Java=Coffee, apa kamu suka minum kopi? Jika ia kamu akan suka Java, Jika tidak kamu akan menemukan cintamu disini.
\end{intermeso}

Java adalah sebuah bahasa pemrograman tingkat tinggi. Pada panduan ini anda
diharapkan sudah menguasi bahasa pemrograman Java terutama OOP pada Java.
Sebab seluruh gim akan ditulis menggunakan bahasa pemrograman Java ini.
\begin{info}
Enjoy the Coffee!
\end{info}

\newpage


\section{LibGDX}
\begin{intermeso}{Kok LibGDX, KOK JAVA?}
Kata Lead tim ZVKT FaSHioN, Refi bertitah : Libgdx framework buaat java dan matkul alrpo 1 dan 2 di Its pakai java harapanya kalau menang (Amin) prestasi kita bisa dikonversi jadi sks dan bisa jadi skem prestasi.
\end{intermeso}

\begin{info}
libGDX is a cross-platform Java game development framework based on OpenGL (ES) that works on Windows, Linux, macOS, Android, your browser and iOS. 

{\raggedleft
Source :  \url{https://libgdx.com/}.\par
}
\end{info}

\subsection{Apa Itu LibGDX?}
Sejatinya Libgdx adalah sebuah \textit{framework} java yang digunakan 
untuk membuat aplikasi berbasis gui, atau dalam hal ini digunakan untuk 
membuat dan mengembangkan sebuah permainan.

\subsection{Game Yang Dibuat Menggunakan LibGDX}
\begin{enumerate}[itemsep=0pt]
    \item \href{https://store.steampowered.com/app/1127400/Mindustry/}{Mindustry} By AnukenDev.
    \item \href{https://store.steampowered.com/app/979110/Space_Haven/}{Space haven} By Bugbyte Ltd.
    \item \href{https://store.steampowered.com/app/249630/Delver/}{delver} by Priority Interrupt.
    \item \href{https://store.steampowered.com/app/546430/Pathway/}{Pathway} by RobotalityPermalink.
    \item \href{https://rockbitegames.com/games/sandship}{Sandship} by Rockbite Games.
\end{enumerate}

\newpage

\section{Setup LibGDX}
\newpage

\section{Proyek Pertamaku!}
\newpage
\centering
\vspace*{\fill}
\begin{info}
Pada bagian selanjutnya kita akan memulai membuat sebuah game!!!
\end{info}
\vspace*{\fill}
\newpage
\begin{thebibliography}{9}

\bibitem{libgdx}
LibGDX.
\textit{LibGDX Official Website}.
\url{https://libgdx.com/}

\bibitem{java}
Oracle.
\textit{Java Documentation}.
\url{https://docs.oracle.com/en/java/}

\end{thebibliography}
\end{document}